% As referências estão embutidas abaixo: o arquivo referencias.bib é recriado a cada compilação.
\begin{filecontents*}[overwrite]{referencias.bib}
@article{avis2018,
  author  = {Avis, Eric and Ferraz, Claudio and Finan, Frederico},
  title   = {Do government audits reduce corruption? {E}stimating the impacts of exposing corrupt politicians},
  journal = {Journal of Political Economy},
  volume  = {126},
  number  = {5},
  pages   = {1912--1964},
  year    = {2018}
}

@article{boas2014,
  author  = {Boas, Taylor C. and Hidalgo, F. Daniel and Richardson, Neal P.},
  title   = {The spoils of victory: campaign donations and government contracts in {B}razil},
  journal = {The Journal of Politics},
  volume  = {76},
  number  = {2},
  pages   = {415--429},
  year    = {2014}
}

@misc{lei9504,
  organization = {Brasil},
  title        = {Lei n. 9.504, de 30 de setembro de 1997},
  note         = {Estabelece normas para as eleições. Diário Oficial da União, Brasília, DF, 1 out. 1997},
  year         = {1997}
}

@misc{lei12527,
  organization = {Brasil},
  title        = {Lei n. 12.527, de 18 de novembro de 2011},
  note         = {Regula o acesso a informações. Diário Oficial da União, Brasília, DF, 18 nov. 2011},
  year         = {2011}
}

@misc{stf2015,
  organization = {Brasil. Supremo Tribunal Federal},
  org-short    = {Brasil},
  title        = {Ação Direta de Inconstitucionalidade n. 4.650/DF},
  note         = {Relator: Min. Luiz Fux. Julgamento em 17 set. 2015},
  address      = {Brasília, DF},
  year         = {2015}
}

@article{brollotroiano2016,
  author  = {Brollo, Fernanda and Troiano, Ugo},
  title   = {What happens when a woman wins an election? {E}vidence from close races in {B}razil},
  journal = {Journal of Development Economics},
  volume  = {122},
  pages   = {28--45},
  year    = {2016}
}

@article{cavalcanti2018,
  author  = {Cavalcanti, Francisco and Daniele, Gianmarco and Galletta, Sergio},
  title   = {Popularity shocks and political selection},
  journal = {Journal of Public Economics},
  volume  = {165},
  pages   = {201--216},
  year    = {2018}
}

@article{claessens2008,
  author  = {Claessens, Stijn and Feijen, Erik and Laeven, Luc},
  title   = {Political connections and preferential access to finance: the role of campaign contributions},
  journal = {Journal of Financial Economics},
  volume  = {88},
  number  = {3},
  pages   = {554--580},
  year    = {2008}
}

@article{colonnelliprem2022,
  author  = {Colonnelli, Emanuele and Prem, Mounu},
  title   = {Corruption and firms},
  journal = {The Review of Economic Studies},
  volume  = {89},
  number  = {2},
  pages   = {695--732},
  year    = {2022}
}

@article{ferrazfinan2008,
  author  = {Ferraz, Claudio and Finan, Frederico},
  title   = {Exposing corrupt politicians: the effects of {B}razil's publicly released audits on electoral outcomes},
  journal = {The Quarterly Journal of Economics},
  volume  = {123},
  number  = {2},
  pages   = {703--745},
  year    = {2008}
}

@article{ferrazfinan2011,
  author  = {Ferraz, Claudio and Finan, Frederico},
  title   = {Electoral accountability and corruption: evidence from the audits of local governments},
  journal = {American Economic Review},
  volume  = {101},
  number  = {4},
  pages   = {1274--1311},
  year    = {2011}
}

@article{mancuso2015,
  author  = {Mancuso, Wagner Pralon},
  title   = {Investimento eleitoral no {B}rasil: balanço da literatura (2001--2012) e agenda de pesquisa},
  journal = {Revista de Sociologia e Política},
  volume  = {23},
  number  = {54},
  pages   = {155--183},
  year    = {2015}
}

@book{stockwatson2019,
  author    = {Stock, James H. and Watson, Mark W.},
  title     = {Introduction to econometrics},
  edition   = {4},
  address   = {New York},
  publisher = {Pearson},
  year      = {2019}
}
\end{filecontents*}

% =============================================================
% Trabalho final de Econometria I (FGV EESP, 2026-2)
% Classe abnTeX2 (artigo, ABNT NBR 6022). Organização das seções, tabelas e figuras
% no padrão de Ferraz e Finan (2011, AER).
% Compilar: pdflatex -> bibtex -> pdflatex -> pdflatex (ou subir a pasta no Overleaf).
% As tabelas (tabelas/) e figuras (figuras/) são geradas pelos scripts em R.
% =============================================================
\documentclass[article,12pt,oneside,a4paper,english,brazil]{abntex2}

\usepackage[T1]{fontenc}
\usepackage[utf8]{inputenc}
\usepackage{mathptmx}
\usepackage{indentfirst}
\usepackage{microtype}
\usepackage{graphicx}
\usepackage{booktabs}
\usepackage{amsmath}
\usepackage{xcolor}
\usepackage{float}
\usepackage{placeins}
\usepackage[alf]{abntex2cite}

\hypersetup{colorlinks=true, linkcolor=black, citecolor=black, urlcolor=blue,
  pdftitle={Auditorias da CGU e doações de campanha}, pdfauthor={Antonio Gama e Leonardo Sørensen}}

% títulos de seção em Times: seção em negrito, subseção em itálico
\renewcommand{\ABNTEXchapterfont}{\rmfamily\bfseries}
\renewcommand{\ABNTEXsectionfontsize}{\normalsize}
\renewcommand{\ABNTEXsubsectionfont}{\rmfamily\itshape}
\renewcommand{\ABNTEXsubsectionfontsize}{\normalsize}

% legendas: "Tabela 1 – Título" em versalete, centralizada
\captionnamefont{\small\scshape}
\captiontitlefont{\small\scshape}
\captionstyle{\centering}

% campos que a dupla ainda precisa preencher (aparecem em vermelho)
\newcommand{\preencher}[1]{\textcolor{red}{[#1]}}

% notas abaixo de tabelas e figuras
\newcommand{\notas}[1]{\par\vspace{4pt}\parbox{\linewidth}{\footnotesize\textit{Notas}: #1}}

\setlength{\parindent}{1.25cm}
\setlength{\parskip}{0pt}
\setlength{\abovecaptionskip}{4pt}
\setlength{\textfloatsep}{12pt plus 2pt minus 4pt}
\setlength{\intextsep}{10pt plus 2pt minus 2pt}

\begin{document}
\selectlanguage{brazil}
\frenchspacing
\textual
\pagestyle{simple}
\SingleSpacing

% ---------------------------------------------------------------
\begin{center}
{\Large\bfseries As empresas abandonam o partido do prefeito exposto?\\[2pt]
Auditorias da CGU e doações de campanha em 2008 e 2012}\\[14pt]
\renewcommand{\thefootnote}{*}
{\itshape Por} {\scshape Antonio Gama e Leonardo Sørensen}\footnote{FGV EESP. Trabalho final da
disciplina Econometria I (2026-2). Código, dados e declaração de uso de inteligência artificial estão no
Apêndice A.}
\end{center}
\renewcommand{\thefootnote}{\arabic{footnote}}
\setcounter{footnote}{0}

\vspace{2pt}
\begin{quote}
\small\noindent\textbf{Resumo:} \textit{Eleitores punem prefeitos cuja corrupção é revelada antes da
eleição. Este trabalho pergunta se as empresas doadoras fazem o mesmo. Usamos os sorteios do programa de
fiscalização da CGU como experimento natural: entre 678 municípios auditados pela primeira vez, comparamos
aqueles cujo relatório foi publicado antes das eleições de 2008 ou 2012 com aqueles cujo relatório, sobre
a mesma gestão, só saiu depois. A variável dependente é a parcela das doações de pessoas jurídicas no
município que vai para o partido do prefeito. Não encontramos efeito da divulgação: a parcela é
praticamente a mesma nos dois grupos e não cai onde a auditoria encontrou mais irregularidades. O
resultado se repete em modelos de resposta binária, em um exercício de variável instrumental e em testes
de robustez, e não parece vir de informação envelhecida, de doações pouco visíveis ou de reações opostas
de empresas diferentes. A evidência sugere que o doador empresarial, ao contrário do eleitor, não pune a
corrupção revelada, embora a amostra não permita descartar efeitos moderados.}

\smallskip
\noindent\textbf{Palavras-chave:} corrupção; auditorias; financiamento de campanha; eleições municipais.

\medskip
\noindent\textbf{Abstract:} \textit{\foreignlanguage{english}{Voters punish mayors whose corruption is
exposed before an election. We ask whether corporate donors do the same. Using Brazil's randomized
municipal audits, we compare 678 municipalities whose first audit report was released before the 2008 or
2012 elections with those whose report, on the same administration, was released only afterwards. We find
no effect of disclosure on the share of corporate donations going to the mayor's party, regardless of how
much corruption the audit uncovered. Corporate donors, unlike voters, do not seem to punish exposed
corruption.}}

\smallskip
\noindent\textbf{Keywords:} \foreignlanguage{english}{corruption; audits; campaign finance; local elections.}
\end{quote}

% ---------------------------------------------------------------
\section{Introdução}

Até 2015, quando o Supremo Tribunal Federal as declarou inconstitucionais \cite{stf2015}, as doações de
empresas eram uma das principais fontes de financiamento das campanhas brasileiras. A literatura mostra que
esse dinheiro não era desinteressado. \citeonline{boas2014} estimam que empreiteiras que doaram a deputados
federais eleitos do partido governista receberam em contratos públicos pelo menos 14 vezes o valor doado, e
\citeonline{claessens2008} mostram que empresas doadoras de candidatos vitoriosos tiveram mais acesso a
crédito bancário. \citeonline{mancuso2015} resume essa literatura sob o rótulo de ``investimento
eleitoral'': a doação é uma aposta em quem vai controlar o orçamento.

Se a doação é um investimento, o que acontece com ela quando o político financiado é exposto como corrupto?
Sobre o eleitor, a resposta é conhecida. \citeonline{ferrazfinan2008} exploram o programa de fiscalização
por sorteios da Controladoria-Geral da União (CGU) e mostram que a divulgação do relatório de auditoria
antes da eleição de 2004 reduziu a chance de reeleição dos prefeitos com mais irregularidades. Sobre o
doador, não. Nos estudos que usam essas auditorias, as doações de campanha aparecem como característica
do prefeito \cite{brollotroiano2016}, e não como resposta à revelação de corrupção.

Este trabalho pergunta: \textit{quando a corrupção de uma gestão municipal vem a público antes da eleição,
as empresas reduzem as doações ao partido do prefeito?} Para responder, seguimos o desenho de
\citeonline{ferrazfinan2008}. Todos os municípios da amostra foram sorteados e auditados pela CGU; o que
varia é o momento da divulgação. Comparamos municípios cujo relatório foi publicado antes das eleições de
2008 ou 2012 com municípios cujo relatório, referente à mesma gestão, só foi publicado depois. Como a ordem
dos sorteios é aleatória, os dois grupos são comparáveis, e a diferença nas doações pode ser atribuída à
divulgação. Combinamos as medidas de corrupção de \citeonline{avis2018} com as prestações de contas
eleitorais do Tribunal Superior Eleitoral (TSE) em uma base \textit{cross-section} de 678 municípios, cada
um observado uma única vez.

O resultado é um efeito nulo. A parcela das doações empresariais que vai para o partido do prefeito é de
27,6\% nos municípios de controle e de 27,5\% nos tratados. A diferença não muda com controles, não depende
do nível de corrupção revelado e não aparece na probabilidade de receber doação nem no valor doado. O
intervalo de confiança descarta quedas superiores a 6,4 pontos percentuais.

Um efeito nulo admite várias leituras, e testamos as que os dados permitem. O resultado não muda quando
olhamos apenas para relatórios recentes, para a eleição em que o nome do doador já era público durante a
campanha, para os municípios em que o próprio prefeito é candidato ou para construtoras e demais empresas
separadamente. O resultado sugere que a disciplina exercida pelo eleitor não é reforçada por quem financia
as campanhas.

O restante do trabalho está organizado como segue. A Seção~\ref{sec:teoria} discute por que a divulgação
afetaria as doações. A Seção~\ref{sec:contexto} descreve as regras de financiamento e o programa de
sorteios. A Seção~\ref{sec:dados} apresenta os dados e a Seção~\ref{sec:metodo}, a estratégia empírica. Os
resultados estão na Seção~\ref{sec:result}, seguidos dos testes de explicações alternativas na
Seção~\ref{sec:alt}. A Seção~\ref{sec:concl} conclui.

% ---------------------------------------------------------------
\section{Arcabouço conceitual}\label{sec:teoria}

Considere uma empresa que decide quanto doar ao partido do prefeito. O retorno esperado da doação é a
probabilidade de esse grupo continuar no poder multiplicada pelo valor do acesso que a doação compra, como
contratos e pagamentos em dia, menos os custos de estar associada a ele. A divulgação de um relatório
desfavorável antes da eleição pode reduzir esse retorno por três canais, que chamamos em conjunto de
\textit{afastamento}.

O primeiro é o \textit{cálculo eleitoral}. Se o eleitor pune a corrupção revelada
\cite{ferrazfinan2008}, o partido do prefeito fica menos competitivo, e a empresa desloca a aposta para os
adversários. O segundo é a \textit{reputação}: aparecer como financiadora de um político exposto tem custo
perante clientes e eleitores. Esse canal depende de a doação ser visível, e a visibilidade mudou entre as
duas eleições que estudamos (Seção~\ref{sec:contexto}): em 2008, o nome do doador só era conhecido depois
da votação; em 2012, passou a ser divulgado durante a campanha. O terceiro é o \textit{risco legal}. O
relatório é enviado ao Ministério Público e ao Tribunal de Contas, e \citeonline{avis2018} concluem que o
principal custo de uma auditoria para o político não é eleitoral, mas o das punições que ela pode gerar.
Para a empresa que contrata com a prefeitura, a doação documenta um vínculo com uma gestão sob
investigação. Diferentemente da reputação, esse canal não depende de publicidade, porque os órgãos de
controle têm acesso às prestações de contas de qualquer forma.

Na direção contrária está a \textit{lealdade}. Se a doação é parte de uma relação de troca
\cite{boas2014}, a empresa que depende do prefeito tem interesse em mantê-lo no cargo justamente quando
ele está ameaçado, e pode até reforçar o apoio. Esse motivo deve ser mais forte quando o próprio prefeito
é candidato e para empresas de setores que dependem de contratos públicos, como a construção, que são
também os mais afetados pelas auditorias \cite{colonnelliprem2022}.

Os canais têm previsões distintas. O afastamento prevê queda das doações com a divulgação, tanto maior
quanto mais grave o que foi encontrado; a lealdade prevê o contrário. A reputação prevê efeito mais
negativo em 2012 do que em 2008. O cálculo eleitoral e a reputação recaem sobre o candidato, e por isso
devem pesar mais quando o prefeito disputa a reeleição do que quando está impedido pelo limite de dois
mandatos \cite{ferrazfinan2011} ou é substituído pelo partido, que ajusta suas candidaturas à revelação de
corrupção \cite{cavalcanti2018}. O risco legal recai sobre a gestão e não depende de quem é o candidato. Por fim, se
afastamento e lealdade atuam sobre empresas diferentes, o efeito médio pode ser nulo mesmo com reações
fortes nos dois sentidos. O sinal do efeito é, portanto, uma questão empírica, e essas previsões orientam
os testes da Seção~\ref{sec:alt}.

% ---------------------------------------------------------------
\section{Contexto institucional}\label{sec:contexto}

\subsection{Doações de empresas a campanhas municipais}

Nas eleições de 2008 e 2012, pessoas jurídicas podiam doar até 2\% do faturamento bruto do ano anterior a
candidatos, comitês financeiros e diretórios partidários \cite{lei9504}. Toda receita devia constar da
prestação de contas entregue à Justiça Eleitoral até trinta dias depois da eleição, com nome e CNPJ do
doador. Durante a campanha, a lei exigia apenas dois relatórios parciais, em 6 de agosto e 6 de setembro,
e dispensava expressamente a identificação dos doadores neles. Em 2008, portanto, o eleitor só sabia quem
financiou cada candidato depois de votar. Em agosto de 2012, com base na Lei de Acesso à Informação
\cite{lei12527}, o TSE determinou pela primeira vez a divulgação dos doadores já nas prestações parciais.
Nos dois anos, parte do dinheiro chegava aos candidatos por meio de diretórios estaduais e nacionais, o
que impedia ligar o doador original ao candidato final. Essas doações indiretas não entram na nossa medida.

\subsection{O programa de fiscalização por sorteios}

Desde 2003, a CGU sorteia municípios com até 500 mil habitantes para auditar o uso de recursos federais.
A partir de 2004, cada sorteio seleciona cerca de 60 municípios, com probabilidade igual dentro de cada
estado. Os auditores passam uma a duas semanas no município, e o relatório com todas as irregularidades
encontradas é publicado na internet e enviado ao Tribunal de Contas da União, ao Ministério Público e à
Câmara Municipal \cite{avis2018}.

\begin{figure}[!htb]
\centering
\caption{Linha do tempo das eleições e da publicação dos relatórios de auditoria}\label{fig:tempo}
\includegraphics[width=0.92\linewidth]{figuras/figura1.pdf}
\notas{As barras indicam o período em que foram publicados os relatórios de cada grupo de sorteios. São
tratados os municípios com relatório publicado antes da eleição em que a gestão auditada foi julgada
(barras pretas) e controles os municípios com relatório publicado depois dela (barras cinza).
\textit{Fonte}: elaboração própria com dados da CGU.}
\end{figure}

O relatório só pode influenciar os doadores se já for público durante a campanha. A CGU divulga os
relatórios em lote, por sorteio,\footnote{Para datar a publicação, lemos a data de geração de 51 relatórios
no sistema e-Aud da CGU, três por sorteio, do 22º ao 38º. Dentro de cada sorteio, as datas coincidem ou
diferem em no máximo duas semanas.} e a Figura~\ref{fig:tempo} mostra o calendário que resulta disso. Para a eleição de 5 de outubro de
2008, os relatórios dos sorteios 22 a 25 saíram até julho de 2008, e os dos sorteios 26 a 28, a partir de
dezembro. Para a eleição de 7 de outubro de 2012, os relatórios dos sorteios 31 a 35 saíram até março de
2012, e os dos sorteios 36 a 38, a partir de dezembro. Nos dois casos, as auditorias do segundo grupo
ainda examinam essencialmente recursos da gestão que disputou a eleição, mas o resultado só se tornou
público pelo menos dois meses depois dela.


% ---------------------------------------------------------------
\FloatBarrier
\section{Dados}\label{sec:dados}

\subsection{Auditorias e medida de corrupção}

Usamos os arquivos de replicação de \citeonline{avis2018}, que codificam as irregularidades de 1.020
auditorias dos sorteios 22 a 38. Como os autores, tratamos como corrupção as irregularidades classificadas
pelos auditores como médias ou graves; nossa medida é o número delas por ordem de serviço auditada. Esses
arquivos são anonimizados e não identificam o município. Recuperamos o município de cada auditoria
cruzando-a com a lista de municípios sorteados pela CGU e validamos 1.019 das 1.020 auditorias; a que
falhou foi descartada.\footnote{Os arquivos trazem o sorteio, um código de estado e características do
Censo de 2000. Dentro de cada combinação de sorteio e estado, cada auditoria foi pareada a um município
sorteado, sem repetição, minimizando a distância em população, índice de Gini e participação da população
urbana. Como a participação urbana vem do mesmo Censo nas duas fontes, o par correto tem valor idêntico, o
que valida o pareamento.}

\subsection{Doações e partido do prefeito}

O partido do prefeito em cada eleição é o do candidato eleito na eleição anterior (2004 ou 2008), segundo
os resultados do TSE.\footnote{Onde houve eleição suplementar durante o mandato (17 municípios), usamos o
eleito nela. Siglas que mudaram entre 2004 e 2008 foram harmonizadas: PFL para DEM; PL e PRONA para PR;
PAN para PTB.} Das prestações de contas de 2008 e 2012, usamos apenas receitas do tipo ``recursos de pessoas
jurídicas''. Somamos, por município e partido, as doações recebidas pelos candidatos a prefeito, pelos
comitês financeiros municipais (de prefeito e únicos) e, em 2012, pelos diretórios municipais. Repasses de
comitês e partidos a candidatos são excluídos, porque o mesmo dinheiro já aparece como receita original de
quem o recebeu da empresa.

A variável dependente principal é a \textit{parcela} das doações de pessoas jurídicas no município que vai
para o partido do prefeito. Ela mede a escolha do doador entre o grupo no poder e seus adversários e, por
ser uma proporção, não é afetada pelo tamanho do município nem pelo encarecimento das campanhas entre 2008
e 2012. Está definida para os 583 municípios com alguma doação empresarial. Duas variáveis secundárias
cobrem toda a amostra: uma \textit{dummy} igual a um se o partido do prefeito recebeu alguma doação e o
logaritmo de um mais o valor recebido. Medimos as doações no nível do partido no município, e não do
candidato, para não condicionar a amostra à decisão de concorrer, que pode ser afetada pela auditoria.

\FloatBarrier
\subsection{Amostra}

Cada município aparece uma única vez. Das 900 auditorias dos sorteios 22 a 28 e 31 a 38, ficamos com as
que têm município identificado e corrupção medida (862) e, entre elas, apenas com a \textit{primeira}
auditoria de cada município (678), porque quem já foi auditado muda de comportamento \cite{avis2018}. A
regra foi conferida com a lista completa de sorteios da CGU desde 2003. Os sorteios 29 e 30 ficam fora
porque cobrem recursos de duas gestões. A Tabela~\ref{tab:amostra} mostra a composição: 433 municípios
tratados e 245 controles.

\begin{table}[!htb]
\centering
\caption{Sorteios, datas de divulgação e amostra}\label{tab:amostra}
\footnotesize
\input{tabelas/tabela1}
\notas{A coluna (2) traz a data do primeiro relatório do sorteio entre os três consultados no sistema
e-Aud. A coluna (5) conta os municípios auditados pela primeira vez, com município identificado e
irregularidades codificadas; a coluna (6), os que tiveram alguma doação de pessoa jurídica na eleição
indicada. No 36º sorteio, apenas 24 dos 60 municípios têm irregularidades codificadas. \textit{Fonte}:
elaboração própria com dados de CGU, Avis, Ferraz e Finan (2018) e TSE.}
\end{table}

\FloatBarrier
\subsection{Estatísticas descritivas}\label{sec:desc}

A Tabela~\ref{tab:desc} compara os dois grupos. O Painel A antecipa o resultado. O partido do prefeito
recebe 27,6\% das doações empresariais do município entre os controles e 27,5\% entre os tratados, e a
proporção de municípios em que ele recebe alguma doação também é praticamente a mesma (37,1\% e 37,4\%).
Os valores doados são maiores entre os tratados, mas a diferença é imprecisa. O próprio prefeito é
candidato de novo em 52,2\% dos municípios de controle e em 50,1\% dos tratados.

\begin{table}[!htb]
\centering
\caption{Estatísticas descritivas e balanceamento}\label{tab:desc}
\footnotesize
\input{tabelas/tabela2}
\notas{As colunas (1) e (2) trazem as médias de cada grupo. A coluna (3) é a diferença entre elas, obtida
pela regressão de cada variável na \textit{dummy} de tratamento, e a coluna (4), o erro padrão robusto à
heterocedasticidade (HC1). A parcela do partido do prefeito é definida para os 583 municípios com alguma
doação de pessoa jurídica (208 controles e 375 tratados). Características municipais do Censo de 2000,
exceto o PIB, de 2005. *** $p<0{,}01$; ** $p<0{,}05$; * $p<0{,}10$. \textit{Fonte}: elaboração própria com
dados de CGU, Avis, Ferraz e Finan (2018), TSE e IBGE.}
\end{table}

O Painel C mostra que as características municipais são muito parecidas. Nenhuma diferença é
significativa, e em uma regressão da \textit{dummy} de tratamento sobre as seis características e uma
\textit{dummy} de eleição não rejeitamos que todos os coeficientes sejam zero ($F=0{,}37$; $p=0{,}90$). É o
que se espera de um sorteio.

O Painel B pede uma ressalva. O número de irregularidades por auditoria é igual nos dois grupos, mas o
número de ordens de serviço é menor entre os controles, de modo que a medida de corrupção por ordem de
serviço é maior. A diferença vem do escopo das auditorias, que diminuiu com o tempo: a média de ordens de
serviço cai de cerca de 26 para 17 a partir do 34º sorteio, e os controles de 2012 são todos posteriores a
ele. Trata-se de uma mudança na régua, não nos municípios. Por isso, a Seção~\ref{sec:rob} repete as
estimações com a corrupção medida em relação à média do próprio sorteio.

% ---------------------------------------------------------------
\FloatBarrier
\section{Estratégia empírica}\label{sec:metodo}

\subsection{Especificação}

Estimamos por mínimos quadrados ordinários (MQO) a regressão
\begin{equation}\label{eq:mqo}
Y_i=\beta_0+\beta_1\,\mathit{Antes}_i+\beta_2\,\mathit{Corr}_i+\beta_3\,(\mathit{Antes}_i\times\mathit{Corr}_i)
+X_i'\gamma+\delta\,E^{2012}_i+\varepsilon_i ,
\end{equation}
em que $Y_i$ é a parcela das doações de pessoas jurídicas recebida pelo partido do prefeito do município
$i$; $\mathit{Antes}_i$ indica que o relatório foi publicado antes da eleição; $\mathit{Corr}_i$ é o número
de irregularidades por ordem de serviço, centrado na média amostral; $X_i$ reúne os controles municipais;
e $E^{2012}_i$ indica a eleição de 2012. Como $\mathit{Corr}$ está centrada, $\beta_1$ é o efeito da
divulgação para um município com corrupção média, e $\beta_3$ mede quanto esse efeito muda a cada
irregularidade adicional por ordem de serviço. O afastamento prevê $\beta_1<0$ e $\beta_3<0$; a lealdade, o
contrário. Começamos pela regressão simples de $Y$ em $\mathit{Antes}$, que equivale à diferença de médias,
e acrescentamos os demais termos em sequência.

Os controles são os logaritmos da população, da renda per capita e do PIB per capita, a participação da
população urbana, o índice de Gini e a taxa de analfabetismo adulto. Como o tratamento é sorteado, eles
não são necessários para a identificação. Entram por dois motivos. Municípios maiores e mais ricos têm
mais empresas e campanhas mais caras, de modo que os controles absorvem parte da variação das doações e
podem reduzir o erro padrão de $\hat\beta_1$; e corrigem eventuais desequilíbrios ao acaso entre os grupos.
Todos são medidos antes dos sorteios (Censo de 2000 e PIB de 2005) e, portanto, não podem ter sido
afetados pela auditoria. A \textit{dummy} de eleição é necessária porque a proporção de tratados difere
entre 2008 e 2012, anos com regras de transparência e custos de campanha distintos. Em uma das
especificações incluímos \textit{dummies} de estado, porque o sorteio é feito dentro de cada estado.

\subsection{Hipóteses de identificação}

A hipótese central é a exogeneidade do tratamento, $E[\varepsilon_i\mid \mathit{Antes}_i,X_i]=0$. Ela se
apoia em dois fatos. Os municípios são escolhidos por sorteio, de modo que o sorteio em que cada um caiu
não depende de suas características nem das de seu prefeito. E a CGU divulga os relatórios em lote, por
sorteio, de modo que não há escolha, município a município, sobre publicar antes ou depois da eleição.
Assim, $\mathit{Antes}$ é função apenas do sorteio, e fatores omitidos que afetam as doações, como a força
da economia local ou a popularidade do prefeito, não deveriam estar correlacionados com ele. O teste de
balanceamento da Tabela~\ref{tab:desc} é coerente com isso. Uma qualificação é que a amostra com parcela
definida exclui municípios sem nenhuma doação empresarial; a divulgação, porém, não altera a chance de
haver doação no município (coeficiente de 0,020; erro padrão de 0,027).

A corrupção, ao contrário do momento da divulgação, não é sorteada. Municípios mais corruptos diferem em
instituições, estrutura econômica e qualidade dos políticos, e essas características também afetam as
doações. Há, portanto, viés de variável omitida em $\beta_2$, que não tem interpretação causal. Já
$\beta_3$ compara municípios com o mesmo nível de corrupção medida, com e sem divulgação. Como
$\mathit{Antes}$ é exógeno em qualquer nível de corrupção, $\beta_3$ identifica como o efeito da divulgação
varia entre municípios mais e menos corruptos, sem que se possa afirmar que é a corrupção em si, e não
algo associado a ela, que gera a variação.

Três escolhas decorrem dos dados. Não incluímos \textit{dummies} de sorteio porque $\mathit{Antes}$ é
função exata do sorteio, e haveria multicolinearidade perfeita; a centralização de $\mathit{Corr}$, por
sua vez, reduz a correlação entre $\mathit{Antes}$ e o termo de interação de 0,81 para $-$0,05. Usamos
erros padrão robustos à heterocedasticidade (HC1) em todas as estimações, porque a variável dependente é
limitada entre zero e um e tem muitos zeros, e o teste de Breusch-Pagan rejeita a homocedasticidade na
maioria das especificações. E, como cada município aparece uma única vez, tratamos as observações como
independentes.

Por fim, o tamanho da amostra limita o que podemos detectar. O erro padrão da diferença de médias é de
3,3 pontos percentuais. Um teste bilateral a 5\% só rejeita a hipótese nula com 80\% de probabilidade
quando o efeito verdadeiro é de pelo menos 2,8 erros padrão, isto é, cerca de 9 pontos percentuais ou um
quarto do desvio padrão da parcela. Efeitos menores que esse podem existir sem que os detectemos.

\subsection{Resposta binária e variável instrumental}\label{sec:binvi}

Para a variável binária, comparamos o modelo de probabilidade linear (MPL), estimado por MQO, com
\textit{probit} e \textit{logit}. O MPL é simples de interpretar, mas impõe efeitos constantes e pode
prever probabilidades fora do intervalo unitário. O \textit{probit} e o \textit{logit} evitam isso ao
passar o índice linear por uma função de distribuição acumulada (normal e logística, respectivamente).
Como deixam de ser lineares nos parâmetros, não são estimados por MQO, mas por máxima verossimilhança: os
coeficientes escolhidos são os que tornam mais provável a amostra observada. Seus coeficientes não têm
leitura direta, e por isso reportamos os efeitos como diferenças de probabilidade prevista para um
município típico \cite{stockwatson2019}.

O MQO compara relatórios publicados antes e depois da eleição. Há, porém, dois sorteios (26 e 36)
realizados poucos meses antes da eleição cujos relatórios só saíram depois. Para eles, ``sorteado antes'' e
``publicado antes'' não coincidem, como em um experimento com cumprimento imperfeito. Isso permite um
exercício de variável instrumental: a variável endógena é $\mathit{Antes}$, e os instrumentos são $Z_1$
(sorteado antes da eleição) e $Z_2$ (meses entre o sorteio e a eleição), ambos determinados apenas pelo
calendário dos sorteios. A relevância é verificada pela estatística $F$ do primeiro estágio. A restrição de
exclusão exige que ser sorteado antes só afete as doações por meio da publicação do relatório; ela seria
violada se a simples visita dos auditores já sinalizasse problemas aos doadores. Com dois instrumentos e
uma variável endógena, o teste $J$ de sobreidentificação avalia essa hipótese. Como a publicação é
determinada pelo sorteio, não esperamos que o MQO seja viesado; os mínimos quadrados em dois estágios
(MQ2E) servem aqui como verificação.

% ---------------------------------------------------------------
\FloatBarrier
\section{Resultados}\label{sec:result}

\subsection{Efeito da divulgação sobre a parcela das doações}

A Tabela~\ref{tab:mqo} traz as estimativas da equação~(\ref{eq:mqo}). Na coluna (1), sem controles, a
divulgação do relatório antes da eleição está associada a uma parcela 0,1 ponto percentual menor para o
partido do prefeito, com erro padrão de 3,3 pontos. Os controles da coluna (2) praticamente não alteram
nem o coeficiente nem o erro padrão, como se espera quando o tratamento é independente das
características dos municípios. O intervalo de confiança de 95\% vai de $-$6,4 a $+$6,3 pontos percentuais.
Em relação à média do grupo de controle (27,6\%), isso descarta quedas superiores a 23\% na parcela do
partido do prefeito, mas não descarta efeitos menores.

\begin{table}[!htb]
\centering
\caption{Efeito da divulgação do relatório sobre as doações de empresas (MQO)}\label{tab:mqo}
\footnotesize
\input{tabelas/tabela3}
\notas{Nas colunas (1) a (4), a variável dependente é a parcela das doações de pessoas jurídicas no
município recebida pelo partido do prefeito; nas colunas (5) e (6), o logaritmo de um mais o valor recebido
(em R\$). Corrupção é o número de irregularidades médias ou graves por ordem de serviço, centrado na média
(2,68); o percentil 75 corresponde a 0,70 acima da média. Controles municipais: logaritmos da população,
da renda per capita e do PIB per capita, participação urbana, Gini e taxa de analfabetismo. Erros padrão
robustos à heterocedasticidade (HC1) entre colchetes. *** $p<0{,}01$; ** $p<0{,}05$; * $p<0{,}10$.
\textit{Fonte}: elaboração própria.}
\end{table}

As colunas (3) e (4) acrescentam a corrupção e sua interação com a divulgação. O coeficiente da interação
é positivo (0,012 e 0,021), pequeno e não significativo: cada irregularidade adicional por ordem de
serviço, cerca de 0,8 desvio padrão, muda o efeito da divulgação em um a dois pontos percentuais, com erro
padrão de três. Para um município no percentil 75 de corrupção, o efeito estimado é de $+$0,6 a $+$1,1
ponto, também indistinguível de zero. O teste $F$ conjunto de que $\beta_1=\beta_3=0$, isto é, de que a
divulgação não tem efeito em nenhum nível de corrupção, não é rejeitado ($p=0{,}91$ e $p=0{,}76$). O
coeficiente da corrupção é negativo e não significativo; como discutido, ele não tem leitura causal. As
\textit{dummies} de estado da coluna (4) elevam o $R^2$ de 0,03 para 0,08 sem mudar as conclusões.

A Figura~\ref{fig:quartis} mostra o mesmo resultado sem impor forma funcional, no formato usado por
\citeonline{ferrazfinan2008} para a reeleição. Se a divulgação afastasse os doadores dos prefeitos mais
corruptos, a linha dos tratados deveria se descolar para baixo da linha dos controles nos quartis
superiores de corrupção. Não é o que se vê: as duas linhas praticamente se sobrepõem do segundo ao quarto
quartil, e a maior distância, de cinco pontos percentuais, aparece no quartil \textit{menos} corrupto,
bem dentro dos intervalos de confiança.

\begin{figure}[!htb]
\centering
\caption{Doações ao partido do prefeito por nível de corrupção encontrado na auditoria}\label{fig:quartis}
\includegraphics[width=0.94\linewidth]{figuras/figura2.pdf}
\notas{Médias e intervalos de confiança de 95\% por quartil do número de irregularidades médias ou graves
por ordem de serviço. Painel A: municípios com alguma doação de pessoa jurídica ($n=583$). Painel B: todos
os municípios ($n=678$). \textit{Fonte}: elaboração própria com dados de CGU, Avis, Ferraz e Finan (2018)
e TSE.}
\end{figure}

As colunas (5) e (6) usam o logaritmo do valor doado e incluem os municípios sem doação. O coeficiente de
\textit{Antes} é positivo (0,06 e 0,09), mas o erro padrão é quatro vezes maior. O resultado é o mesmo:
nenhum efeito detectável, na média ou em função da corrupção. Também não há efeito sobre a decisão do
partido do prefeito de lançar candidato (0,047; erro padrão de 0,039).

\FloatBarrier
\subsection{Probabilidade de receber doação}

A Tabela~\ref{tab:bin} analisa a margem extensiva. No MPL, a divulgação muda a probabilidade de o partido do
prefeito receber alguma doação empresarial em 0,3 ponto percentual, com erro padrão de 3,8 pontos. Os
coeficientes do \textit{probit} e do \textit{logit} não são diretamente comparáveis aos do MPL, mas as
probabilidades previstas são. Para um município típico (controles na média, eleição de 2008), a
probabilidade prevista passa de 0,360 para 0,365 no \textit{probit} e de 0,359 para 0,363 no \textit{logit}:
um efeito de 0,4 ponto percentual, praticamente igual ao do MPL. Para comparação, um desvio padrão a mais
no logaritmo da população eleva a mesma probabilidade em 4,6 pontos no \textit{probit}. Os três modelos
coincidem porque as probabilidades previstas ficam longe de zero e de um (entre 0,13 e 0,77 no MPL, sem
nenhuma previsão fora do intervalo unitário), região em que as funções de ligação são aproximadamente
lineares. Incluindo a interação com a corrupção no \textit{probit}, o teste conjunto de que a divulgação
não tem efeito tampouco é rejeitado ($p=0{,}98$).

\begin{table}[!htb]
\centering
\caption{Probabilidade de o partido do prefeito receber doação de empresa}\label{tab:bin}
\footnotesize
\input{tabelas/tabela4}
\notas{A variável dependente é igual a um se o partido do prefeito recebeu alguma doação de pessoa jurídica.
MPL estimado por MQO, com erros padrão robustos (HC1); \textit{probit} e \textit{logit} por máxima
verossimilhança. O município típico tem controles na média amostral, corrupção média e eleição de 2008.
Erros padrão entre colchetes. *** $p<0{,}01$; ** $p<0{,}05$; * $p<0{,}10$. \textit{Fonte}: elaboração
própria.}
\end{table}

\FloatBarrier
\subsection{Variável instrumental}

A Tabela~\ref{tab:vi} apresenta o exercício descrito na Seção~\ref{sec:binvi}. Na amostra com parcela
definida, 55 municípios dos sorteios 26 e 36 foram sorteados antes da eleição e tiveram o relatório
publicado depois. O primeiro estágio é forte: ser sorteado antes eleva em 87 pontos percentuais a
probabilidade de publicação antes da eleição, e a estatística $F$ está muito acima do valor de referência
de 10. Ela é tão alta porque, dado o calendário, a publicação é quase uma função exata dos instrumentos. A
forma reduzida (coluna 2) mostra efeito de 0,9 ponto percentual de ser sorteado antes. O MQ2E com um
instrumento (coluna 3), que equivale à forma reduzida dividida pelo primeiro estágio, estima efeito de 1,1
ponto, com erro padrão de 4,0; com dois instrumentos (coluna 4), $-$0,4 ponto. O teste $J$ não rejeita a
validade conjunta dos instrumentos ($p=0{,}36$). MQO e MQ2E contam, portanto, a mesma história.

\begin{table}[!htb]
\centering
\caption{Variável instrumental: sorteio antes da eleição como instrumento para a divulgação}\label{tab:vi}
\footnotesize
\input{tabelas/tabela5}
\notas{A variável dependente é a parcela das doações de pessoas jurídicas recebida pelo partido do prefeito.
$Z_1$: município sorteado antes da eleição; $Z_2$: meses entre o sorteio e a eleição. A estatística $F$
testa a exclusão dos instrumentos no primeiro estágio (robusta, HC1). Teste $J$ de sobreidentificação com
um grau de liberdade. Erros padrão robustos (HC1) entre colchetes. \textit{Fonte}: elaboração própria.}
\end{table}

\FloatBarrier
\subsection{Robustez}\label{sec:rob}

\begin{table}[!htb]
\centering
\caption{Testes de robustez}\label{tab:rob}
\footnotesize\setlength{\tabcolsep}{4.5pt}
\input{tabelas/tabela6}
\notas{Todas as regressões incluem a medida de corrupção, os controles municipais e \textit{dummies} de
estado; as colunas (1) a (4) incluem a \textit{dummy} de eleição. Na coluna (3), a corrupção é o desvio em
relação à média do sorteio; na coluna (4), o logaritmo do número de irregularidades (centrado), com o
logaritmo das ordens de serviço como controle. Erros padrão robustos (HC1) entre colchetes.
*** $p<0{,}01$; ** $p<0{,}05$; * $p<0{,}10$. \textit{Fonte}: elaboração própria.}
\end{table}

A Tabela~\ref{tab:rob} repete a especificação mais completa (coluna 4 da Tabela~\ref{tab:mqo}) sob cinco
variações. A coluna (2) exclui os sorteios 26 e 36, cujos municípios receberam os auditores antes da
eleição sem que o relatório fosse publicado; o 36º é também o sorteio com auditorias incompletas
(Tabela~\ref{tab:amostra}). As colunas (3) e (4) trocam a medida de corrupção: primeiro, pelo desvio em
relação à média do próprio sorteio, o que elimina a mudança de escopo das auditorias discutida na
Seção~\ref{sec:desc}; depois, pela especificação de \citeonline{avis2018}, com o logaritmo do número de
irregularidades e o logaritmo das ordens de serviço como controle. As colunas (5) e (6) estimam cada
eleição separadamente. Em todos os casos, e para as duas variáveis dependentes, o coeficiente de
\textit{Antes} fica entre $-$1,2 e $+$3,6 pontos percentuais e o da interação é positivo ou próximo de
zero, sempre com erros padrão maiores que as estimativas. Excluir os 17 municípios com eleição suplementar
também não altera o resultado ($-$0,009; erro padrão de 0,033).


% ---------------------------------------------------------------
\FloatBarrier
\section{Explicações alternativas}\label{sec:alt}

Um efeito nulo pode significar que os doadores não reagem à corrupção revelada, mas também pode esconder
uma reação que o desenho não capta. Esta seção examina cinco explicações desse segundo tipo. Para cada uma,
dizemos o que ela prevê, como a testamos e o que os dados mostram. Os testes estão na
Tabela~\ref{tab:alt}, e todas as regressões usam os controles municipais e a \textit{dummy} de eleição.

\begin{table}[!htb]
\centering
\caption{Testes de explicações alternativas}\label{tab:alt}
\footnotesize\setlength{\tabcolsep}{3pt}
\input{tabelas/tabela7}
\notas{Nas colunas (1), (3) e (4), a variável dependente é a parcela das doações de pessoas jurídicas
recebida pelo partido do prefeito; na coluna (2), a parcela recebida pelo partido pelo qual o prefeito
concorre (ou pelo partido da eleição anterior, se ele não concorre). Nas colunas (5) e (6), \textit{dummies}
iguais a um se o partido do prefeito recebeu doação de empresa de construção civil e obras ou de empresa
de outro setor, na eleição de 2012 (MPL). A coluna (1) inclui a \textit{dummy} de prefeito candidato. Erros
padrão robustos (HC1) entre colchetes. \textit{Fonte}: elaboração própria.}
\end{table}

\subsection{O alvo do doador é o prefeito, não o partido}

Medimos as doações ao partido que venceu a eleição
anterior, mas a relação da empresa pode ser com a pessoa do prefeito. Em 51\% dos municípios ele é
candidato de novo, e 88 desses prefeitos concorrem por outra legenda, levando consigo doações que nossa
medida não conta. Se o afastamento recai sobre a pessoa, ele deveria aparecer sobretudo quando o prefeito
está na disputa: é aí que o cálculo eleitoral e a reputação pesam, e também onde a lealdade teria mais a
proteger. A coluna (1) permite que o efeito da divulgação seja diferente nesses municípios. A interação
é negativa ($-$5,0 pontos percentuais), na direção prevista pelo afastamento, mas o erro padrão é de 6,3
pontos. Somando os dois termos, o efeito para prefeitos candidatos é de $-$2,8 pontos, com intervalo de
confiança de 95\% entre $-$12,6 e $+$7,0. A comparação só é válida se a própria decisão de concorrer não
for afetada pela auditoria, e a Tabela~\ref{tab:desc} mostra que não é: a proporção de prefeitos
candidatos difere em apenas 2,1 pontos entre os grupos, sem significância. A coluna (2) ataca o mesmo
problema por outro lado: refaz a variável dependente seguindo o prefeito até o partido pelo qual ele
concorre. O efeito estimado é de $+$1,0 ponto, também nulo.

\subsection{Informação envelhecida}

Os relatórios dos tratados saíram entre 3 e 24 meses antes da eleição.
Se os doadores esquecem, ou se a campanha produz informações mais recentes que se sobrepõem à auditoria, o
efeito deveria concentrar-se nos relatórios publicados perto da votação. A coluna (3) divide os tratados
em dois grupos: relatórios publicados nos doze meses anteriores à eleição (sorteios 23 a 25, 34 e 35) e
relatórios mais antigos. Os coeficientes são de $+$0,5 e $-$1,0 ponto, e não diferem entre si
($p=0{,}71$). Não há sinal de um efeito que se dissipa com o tempo: mesmo os relatórios recentes, que
chegam no auge da pré-campanha, não alteram a parcela do partido do prefeito.

\subsection{Visibilidade da doação}

O custo de reputação exige que o público saiba quem doou. Como vimos na
Seção~\ref{sec:contexto}, em 2008 os nomes dos doadores só eram divulgados depois da eleição, enquanto em
2012 passaram a sair durante a campanha. A mudança de regra funciona como um teste do canal de reputação:
se ele importasse, a divulgação do relatório deveria reduzir mais as doações em 2012 do que em 2008. A
coluna (4) mostra o contrário de uma diferença: o efeito é de $-$0,1 ponto em 2008, e a interação com 2012
é de $+$0,2 ponto, com erro padrão de 6,5. O resultado não descarta que a reputação importe para
empresas de grande porte, mais expostas à imprensa, mas indica que ela não move o doador típico de uma
campanha municipal.

\subsection{Reações opostas que se cancelam}

O arcabouço da Seção~\ref{sec:teoria} prevê que empresas que dependem
de contratos com a prefeitura tenham motivo para manter o apoio, enquanto as demais se afastam. Se as duas
reações tiverem tamanhos parecidos, o efeito médio seria nulo mesmo com reações fortes. As prestações de
contas de 2012 informam o setor do doador, e as empresas de construção civil e obras, as mais dependentes de
contratos municipais, respondem por 21\% do valor doado. As colunas (5) e (6) estimam o efeito da
divulgação separadamente para a probabilidade de o partido do prefeito receber doação de construtoras e de
outras empresas. Os dois coeficientes são negativos e parecidos ($-$3,2 e $-$3,0 pontos) e nenhum é
significativo. O padrão de cancelamento exigiria sinais opostos, o que não aparece. Esta é, porém, a
evidência mais frágil da tabela: só cobre 2012, e apenas 8,6\% dos partidos de prefeito recebem doação de
construtoras, de modo que uma queda de 3,2 pontos seria grande em termos relativos se fosse precisa.

\subsection{Contaminação dos controles}

Os municípios dos sorteios 26 e 36 receberam os auditores antes da
eleição, embora o relatório só tenha saído depois. Se a simples visita já alertasse os doadores, esses
municípios seriam controles parcialmente tratados, e a comparação subestimaria o efeito. Duas evidências vão
contra essa explicação. Excluir esses sorteios não altera a estimativa (coluna 2 da Tabela~\ref{tab:rob}).
E, no exercício de variável instrumental, a forma reduzida, que compara diretamente quem foi sorteado antes
e depois da eleição, também é nula (Tabela~\ref{tab:vi}).

\subsection{Explicações que os dados não permitem testar}

Duas explicações ficam fora do alcance dos dados. A primeira é que o relatório não traga informação nova:
empresas que fornecem à prefeitura acompanham licitações e pagamentos de perto e podem conhecer as
irregularidades antes dos auditores. Essa leitura é coerente com a ausência de efeito mesmo nos municípios
mais corruptos (Figura~\ref{fig:quartis}), mas não pode ser separada de uma simples indiferença do doador.
A segunda é que a reação ocorra em doações que não observamos: a empresa pode manter o apoio por meio de
diretórios estaduais e nacionais, que repassam os recursos sem identificar o doador original, ou fora da
contabilidade oficial. Nesse caso, o que muda é o canal da doação, não o seu destino, e nossa medida não o
captaria.

% ---------------------------------------------------------------
\FloatBarrier
\section{Conclusão}\label{sec:concl}

Este trabalho perguntou se as empresas reduzem as doações ao partido do prefeito quando a corrupção da
gestão é revelada antes da eleição. Usando a ordem aleatória dos sorteios da CGU e as eleições de 2008 e
2012, não encontramos evidência disso. A parcela das doações empresariais destinada ao partido do
prefeito é praticamente idêntica com e sem divulgação prévia do relatório, e o efeito não cresce com a
gravidade do que foi encontrado.

A leitura econômica mais direta é que o doador empresarial não se comporta como o eleitor de
\citeonline{ferrazfinan2008}. Se a doação é um investimento em uma relação já estabelecida, a empresa tem
pouco a aprender com o relatório e pouco a ganhar rompendo com quem ainda controla o orçamento. A
implicação é que a transparência disciplina o político pelo voto e pela Justiça, mas não pelo
financiamento: quem paga a campanha não reforça a punição aplicada por quem vota.

Há limitações. A principal é o poder estatístico: só detectaríamos com segurança efeitos de cerca de nove
pontos percentuais, e efeitos moderados não podem ser descartados. O comitê financeiro único mistura as
campanhas de prefeito e de vereador. A data de publicação é aproximada pela data de geração do arquivo do relatório. A validade
externa é restrita a municípios com até 500 mil habitantes, ao período 2008--2012 e a um regime em que a
doação empresarial era legal. Por fim, o tratamento varia por sorteio; erros
padrão agrupados nesse nível, fora do escopo da disciplina, seriam uma extensão natural.

\bibliography{referencias}

% ---------------------------------------------------------------
\section*{Apêndice A -- Código, dados e uso de inteligência artificial}

\noindent\textbf{Código (R).} Repositório: \preencher{link do repositório}. A pasta \texttt{scripts}
reproduz todas as tabelas e figuras a partir da base final (\textit{scripts} 02 a 07). A pasta
\texttt{construcao\_base} documenta como a base final foi montada a partir das bases de origem
(\texttt{01\_montar\_base.R}) e traz os \textit{scripts} em Python da consolidação inicial (chave de
municípios e pareamento das auditorias).

\smallskip
\noindent\textbf{Base de dados.} A base final (\texttt{base\_final.csv}, uma linha por município) está na
pasta \texttt{dados} do repositório. As bases de origem, identificadas pelo código IBGE (CGU, Avis, Ferraz
e Finan, TSE, Censo e PIB), estão em \preencher{link da pasta com as bases}.

\smallskip
\noindent\textbf{Uso de inteligência artificial.} Usamos o assistente Claude (Anthropic) na discussão do
tema e do desenho empírico, na organização das bases (chave de municípios, pareamento dos arquivos de Avis,
Ferraz e Finan e leitura das datas dos relatórios) e na elaboração dos \textit{scripts} em R e da primeira
versão do texto e das tabelas. As decisões sobre a pergunta, a amostra e as especificações foram tomadas
por nós, que revisamos o código e os resultados e respondemos pelo conteúdo. Conversas: (i) escolha do
tema, desenho e organização das bases: \preencher{link público da conversa 1}; (ii) \textit{scripts},
tabelas e texto: \preencher{link público da conversa 2}.

\end{document}
